# SNAPSHOT_SOURCE: lindicaphxag-tech/cs@0b26a5b0853e3b136a92c17e614a3a1936673621 path=papers/eeg-default-regularisation/paper_vnext/01_intro.tex
\section{Introduction}

Cross-subject EEG decoding asks a model trained on one group of participants
to operate on a new person who contributes no labelled calibration trials.
This setting makes regularisation scientifically consequential rather than
merely numerical. Covariance and tangent-space representations can contain
hundreds or thousands of coordinates, while the number of independent source
participants is modest and inter-participant covariance structure is
heterogeneous. Analytic shrinkage therefore provides an attractive default
\cite{ledoit2004,chen2010oas},
but the statistical objective used to stabilise a covariance matrix is not the
same quantity on which a transferred decoder is evaluated.

For a two-class linear discriminant, let
\[
 A_\lambda=(1-\lambda)S+\lambda\nu I,
 \qquad
 w_\lambda=A_\lambda^{-1}d_s,
\]
where \(S\) and \(d_s\) are source within-class covariance and class
contrast. Covariance-fidelity criteria choose \(\lambda\) according to how
well \(A_\lambda\) predicts covariance structure. A held-participant decoder
instead depends on how the inverse-spectrum weighting in \(w_\lambda\)
aligns with that participant's class contrast. These objectives may coincide,
but established covariance-based EEG decoding, including Riemannian
classifiers and filter-bank spatial patterns \cite{barachant2012,ang2012fbcsp},
already provides a framework in which classification risk and covariance
estimation are distinct targets. Our
question is therefore narrower and empirical: \emph{does this distinction
produce a reproducible validity problem in cross-subject EEG when target
information is unavailable?}

Recent EEG research has explored additional model regularisers,
optimal-transport adaptation, cyclic representation transfer and adaptive
spatial--temporal attention
\cite{miao2025freezing,morvay2026wbt,he2026cyclic,ren2026asta}.
Riemannian spatial filtering has also been evaluated with ten decoding
algorithms on six public MI datasets \cite{pan2025rsf}. These studies
establish a high empirical bar, but their primary questions and target
information budgets differ. In particular, unsupervised target-domain
adaptation cannot be ranked against strict zero-target-information (Z0)
selection as though they used identical participant information. Our
primary contribution is thus a mechanism and measurement-validity analysis
inside the predeclared Z0 covariance/tangent-linear-discriminant family,
not a claim to outperform modern domain-adaptive encoders.

The operator geometry makes the distinction precise. For an evaluation
participant with covariance \(\Sigma_g\) and contrast \(d_g\), define the
standardized margin
\[
 m_g(\lambda)=
 \frac{w_\lambda^\top d_g}
 {\sqrt{w_\lambda^\top\Sigma_gw_\lambda}}.
\]
The regularisation-path derivative \(w_\lambda'\) can be decomposed into a
component that only rescales \(w_\lambda\) and a
\(\Sigma_g\)-orthogonal rotation \(r_g\). The scale component cancels from
the standardized margin, giving
\[
 m_g'(\lambda)=
 \frac{r_g(\lambda)^\top d_g}
 {\sqrt{w_\lambda^\top\Sigma_gw_\lambda}}.
\]
In contrast, the covariance-predictive NLL derivative contains covariance
operators but no class-contrast vector. Hence a covariance optimum is not, in
general, a discriminative stationary point. Around an interior covariance
optimum, covariance regret is second order in path displacement whereas a
non-zero margin derivative is first order. This supplies a mechanistic reason
why a shallow covariance-fidelity basin can admit materially different
decisions without claiming that such fragility must occur in every EEG cohort.

The historical version of this study could not establish that chain. Its
regularisation grid excluded the near-zero analytic-shrinkage region; a
proposed transfer proxy squared a cross-participant margin and therefore lost
direction information; several summary-table gains could not be traced to one
machine-readable configuration lineage; and an apparent filter-bank advantage
mixed representation and regularisation effects. We therefore reconstruct the
study from participant-level artifacts, full-domain paths and explicitly
frozen evidence tiers. Historical rows with unresolved arithmetic or
provenance are quarantined rather than repaired by inference.

The first prospective test, EEGMMIDB S021--S040, was frozen before any target
outcome and used a strict Z0 information budget: no target labels, target
feature moments, target-fitted geometry or target distribution statistics
entered representation or regularisation selection. It tests whether the
source objective separation has a held-participant consequence. Negative
algorithm results are retained with equal status: lower-tail signed selection
failed to improve over correctly nested source tuning, and a
covariance-fidelity-constrained selector was prospectively inferior on a
separate EEGMMIDB block. These failures prevent the paper from turning a
measurement-validity finding into an unsupported new-estimator story.

We next recover a stronger external evidence tier. BNCI2014\_002 and
BNCI2015\_001 were frozen under one TRUE PRE-DATA protocol before project
access to their outcomes. Each cohort uses its native montage, full outer
leave-one-participant-out evaluation, participant-grouped inner selection and a
dense regularisation path whose target scores are frozen before AUC
calculation. These datasets test both source-objective separation and a
post-S021 Regularisation Validity Frontier. The frontier does not ask for a
target-optimal \(\lambda\); it asks how much held-participant AUC variation
is compatible with a specified amount of source covariance-fidelity regret.
A pre-outcome mathematical correction replaces a structurally meaningless
generic zero-crossing statistic with material-fragility budgets for 0.01,
0.02 and 0.05 AUC envelope spans.

A separate prospectively frozen EEGMMIDB S041--S060 experiment asks a
different scientific question: can regularisation objective change the
conclusion drawn about EEG representations? It compares a wide-band
231-dimensional tangent representation, an eight-band 1848-dimensional
filter-bank representation, a source-only matched-dimension filter-bank PCA
control and a label-independent nuisance-dimension control. This design
permits the original hidden-filter-bank hypothesis to fail. A source-budget
intervention further tests whether a simple monotonic feature-to-sample ratio
law is sufficient to explain any objective mismatch.

Accordingly, the contribution is an evidence chain rather than a leaderboard.
We (i) derive the path-rotation condition that separates covariance fidelity
from discriminative transfer; (ii) test the complete scalar domain under
participant-grouped source selection; (iii) prospectively evaluate the target
consequence and its harmed-participant tail; (iv) replicate source-objective
separation in two TRUE PRE-DATA external cohorts while preserving
dataset-level inferential boundaries; (v) quantify corrected, uncertainty-aware
regularisation validity frontiers; and (vi) prospectively test whether the
choice of regularisation objective changes a representation conclusion. We do
not claim a universal EEG regularisation law, a universal multiplier, or a
state-of-the-art replacement covariance estimator.
